\chapter{Richiami di Algebra Lineare, rototraslazioni e formula fondamentale del moto dei corpi rigidi}
\begin{definition}{}
Una matrice $A\in\mathbb{R}^{n\times n}$ si dice \emph{ortogonale} se
\[
A^{\mathsf T}A=I_n
\]
(equivalentemente $AA^{\mathsf T}=I_n$), cioè se le sue colonne (e righe) formano un insieme ortonormale.
\end{definition}

\begin{oss}{}
Se $A$ è ortogonale, allora
\[
\det(A)=\pm 1.
\]
infatti ricordando che $\det(A) = \det(A^T)$ e la formula di Binet $\det(AB) = \det(A) \det(B)$ otteniamo
\[
\det(A^T A) = \det(A) \det(A^T) = \det(A)^2 = \det(I) = 1 \Rightarrow \det(A) = \pm 1
\]
Con la classica identificazione $\operatorname{End}(\mathbb{R}^n) \cong \mathbb{R}^{n \times n }$ tramite $A \to L_A: \mathbb{R}^n \to \mathbb{R}^n , \underline{x} \mapsto A\underline{x}$ chiamiamo le matrici ortogonali con determinante 1 rotazioni e le matrici ortogonali con determinante -1 riflessioni.
\end{oss}

\begin{example}{}
Consideriamo
\[
A=
\begin{pmatrix}
0 & 1 & 0\\
1 & 0 & 0\\
0 & 0 & 1
\end{pmatrix}.
\]
Si ha $A^{\mathsf T}A=I$ e $\det(A)=-1$, quindi $A$ è ortogonale ma non è una rotazione (inverte l'orientazione).

L'applicazione lineare associata $L_A(\underline{x})=A\underline{x}$ scambia le prime due coordinate e lascia invariata la terza:
\[
L_A(x,y,z)=(y,x,z).
\]
Geometricamente, è la riflessione rispetto al piano $x=y$ (equivalentemente: scambia gli assi $x$ e $y$ e lascia fisso l'asse $z$).
\end{example}

\begin{example}{}
Consideriamo
\[
A=
\begin{pmatrix}
1 & 0 & 0\\
0 & 1 & 0\\
0 & 0 & -1
\end{pmatrix}.
\]
Anche in questo caso $A^{\mathsf T}A=I$ e $\det(A)=-1$, quindi $A$ è ortogonale e inverte l'orientazione.

L'applicazione lineare associata $L_A(\underline{x})=A\underline{x}$ lascia invariati $x$ e $y$ e cambia segno a $z$:
\[
L_A(x,y,z)=(x,y,-z).
\]
Geometricamente, è la riflessione rispetto al piano $z=0$ (piano $xy$).
\end{example}
Noi siamo interessati per lo più alle \emph{rotazioni}, ovvero alle matrici A tali che $\det(A) = 1$.

\begin{example}{}
Sia
\[
R=
\begin{pmatrix}
\cos\theta & -\sin\theta & 0\\
\sin\theta & \cos\theta  & 0\\
0          & 0           & 1
\end{pmatrix}.
\]
Si ha $\det R=1$ e
\[
RR^{\mathsf T}
=
\begin{pmatrix}
\cos\theta & -\sin\theta & 0\\
\sin\theta & \cos\theta  & 0\\
0          & 0           & 1
\end{pmatrix}
\begin{pmatrix}
\cos\theta & \sin\theta & 0\\
-\sin\theta& \cos\theta & 0\\
0          & 0          & 1
\end{pmatrix}
=
\begin{pmatrix}
1&0&0\\
0&1&0\\
0&0&1
\end{pmatrix}.
\]
È la rotazione attorno all'asse $z$ di un angolo $\theta$ in senso antiorario.

\begin{center}
\begin{tikzpicture}[scale=1.05, line cap=round, line join=round]
  % assi 3D in prospettiva (lasciati identici ai tuoi)
  \coordinate (O) at (0,0);
  \draw[->,line width=1.1pt] (O) -- (-2.2,-1.3) node[below left] {$x$};
  \draw[->,line width=1.1pt] (O) -- (3.4,0) node[below right] {$y$};
  \draw[->,line width=1.1pt] (O) -- (0,3.2) node[above] {$z$};

  % archetto di rotazione nel piano xy (Sdraiato sul pavimento!)
  % Parte vicino all'asse x (angolo 210) e va verso l'asse y (angolo 340)
  \draw[->,line width=1.0pt, red] (-1.1,-0.65) arc[start angle=210,end angle=340, x radius=1.3, y radius=0.6];

  % Etichetta theta posizionata correttamente
  \node[red] at (0.2, -0.8) {$\theta$};

  % Opzionale ma consigliato: freccia di omega sull'asse z
  \draw[->,line width=1.2pt, blue] (O) -- (0,1.8) node[right] {$\underline{\omega}$};
\end{tikzpicture}
\end{center}


\inizioapprof Calcoliamo autovalori e autospazi di $R$.
\[
\det(R-\lambda I)=
\begin{vmatrix}
\cos\theta-\lambda & -\sin\theta & 0\\
\sin\theta & \cos\theta-\lambda & 0\\
0 & 0 & 1-\lambda
\end{vmatrix}
\]
\[
=(1-\lambda)
\begin{vmatrix}
\cos\theta-\lambda & -\sin\theta\\
\sin\theta & \cos\theta-\lambda
\end{vmatrix}
=(1-\lambda)\bigl[(\cos\theta-\lambda)^2+\sin^2\theta\bigr]
\]
\[
=(1-\lambda)\bigl[\cos^2\theta+\lambda^2-2\cos\theta\,\lambda+\sin^2\theta\bigr]
=(1-\lambda)\bigl[\lambda^2-2\cos\theta\,\lambda+1\bigr]
=p(\lambda).
\]
Quindi $p(\lambda)=0$ se e solo se $\lambda_1=1$ oppure $\lambda^2-2\cos\theta\,\lambda+1=0$.
Il discriminante è
\[
\Delta=4\cos^2\theta-4=4(\cos^2\theta-1)=-4\sin^2\theta=i^2\,4\sin^2\theta,
\]
e dunque
\[
\lambda_{2,3}=\frac{2\cos\theta\pm 2i|\sin\theta|}{2}
=\cos\theta\pm i|\sin\theta|.
\]

Concludiamo quindi che:
\begin{itemize}
\item se $\theta=2k\pi$ allora $\lambda_1=\lambda_2=\lambda_3=1$;
\item se $\theta=(2k+1)\pi$ allora $\lambda_1=1$ e $\lambda_2=\lambda_3=-1$;
\item mentre se $\theta\neq k\pi$ allora $\lambda_1=1$ è l'unico autovalore reale.
\end{itemize}

Calcoliamo gli autospazi.
\begin{itemize}
\item $\theta\neq k\pi \ \Rightarrow\ E_1(R)=\Ker(R-I)$ con
\[
R-I=
\begin{pmatrix}
\cos\theta-1 & -\sin\theta & 0\\
\sin\theta & \cos\theta-1 & 0\\
0 & 0 & 0
\end{pmatrix},
\]
da cui il sistema
\[
(\cos\theta-1)x-\sin\theta\,y=0,\qquad y=0,\qquad z\in\mathbb{R},
\]
e quindi
\[
E_1(R)=\mathrm{span}\!\left\{\begin{pmatrix}0\\0\\1\end{pmatrix}\right\}.
\]
Ovvero l'unico sottospazio che viene mandato in se stesso è l'asse $z$, che è l'asse di rotazione.

\item $\theta=2k\pi \ \Rightarrow\ E_1(R)=\Ker(R-I)$ con $R-I=0$, quindi
\[
E_1(R)=\mathbb{R}^3.
\]
In questo caso non c'è rotazione e tutto lo spazio viene mandato in se stesso.

\item $\theta=(2k+1)\pi \ \Rightarrow\ E_1(R)=\Ker(R-I)$ e $E_{-1}(R)=\Ker(R+I)$.
In particolare
\[
R-I=
\begin{pmatrix}
-2&0&0\\
0&-2&0\\
0&0&0
\end{pmatrix}
\quad\Rightarrow\quad
E_1(R)=\mathrm{span}\!\left\{\begin{pmatrix}0\\0\\1\end{pmatrix}\right\},
\]
e
\[
R+I=
\begin{pmatrix}
0&0&0\\
0&0&0\\
0&0&-2
\end{pmatrix}
\quad\Rightarrow\quad
E_{-1}(R)=\mathrm{span}\!\left\{\begin{pmatrix}1\\0\\0\end{pmatrix},
\begin{pmatrix}0\\1\\0\end{pmatrix}\right\}.
\]
Otteniamo quindi che $R$ in questo caso mantiene fisso l'asse $z$ (asse di rotazione) e riflette i vettori sul piano $xy$ lungo la loro direzione.
\end{itemize}

\begin{oss}{}
Da $R$ possiamo ottenere le rotazioni attorno agli assi $x$ e $y$ semplicemente moltiplicando a destra $R$ per le matrici
\[
T_x=
\begin{pmatrix}
0&1&0\\
0&0&1\\
1&0&0
\end{pmatrix},
\qquad
T_y=
\begin{pmatrix}
0&0&1\\
1&0&0\\
0&1&0
\end{pmatrix}.
\]
\end{oss}
\fineapprof
\begin{oss}{}
Essendo $R^{-1}=R^{\mathsf T}$ allora $R^{\mathsf T}$ è una rotazione attorno all'asse $z$ di un angolo $\theta$ in senso orario.
\end{oss}
\end{example}

\begin{definition}{}
Sia $A\in\mathbb{R}^{n\times n}$.
\begin{itemize}
\item $A$ si dice \emph{simmetrica} se $A^{\mathsf T}=A$.
\item $A$ si dice \emph{antisimmetrica} se $A^{\mathsf T}=-A$.
\end{itemize}
\end{definition}

Applichiamo i richiami di Algebra Lineare al corpo rigido.\\
\inizioapprof
Per caratterizzare la configurazione del corpo rigido ne scelgo un punto che a $t=0$ avrà posizione $\underline{x}_0$.

Nel tempo il punto si muove e il corpo rigido ruota attorno al punto:
\begin{enumerate}[label=\roman*)]
\item $\underline{x}_p(t)$;
\item fisso una matrice di rotazione attorno a $\underline{x}_p(t)$.
\end{enumerate}
Notiamo che servono 6 informazioni: le 3 coordinate di $\underline{x}_p(t)$ e 3 coefficienti per $R$.
\\Una matrice generica $3\times 3$ ha $9$ coefficienti liberi. Una matrice di rotazione deve invece soddisfare i vincoli
\[
R^{\mathsf T}R=I,\qquad \det R=1,
\]
cioè appartenere a $SO(3)$.

Scrivendo $R=(\underline r_1\,|\,\underline r_2\,|\,\underline r_3)$, la condizione $R^{\mathsf T}R=I$ equivale a dire che le colonne sono ortonormali:
\[
\langle \underline r_i,\underline r_j\rangle=\delta_{ij}.
\]
Questa condizione impone $6$ vincoli indipendenti:
\begin{itemize}
\item $3$ vincoli di norma: $\|\underline r_1\|=1,\ \|\underline r_2\|=1,\ \|\underline r_3\|=1$;
\item $3$ vincoli di ortogonalità: $\underline r_1\cdot \underline r_2=0,\ \underline r_1\cdot \underline r_3=0,\ \underline r_2\cdot \underline r_3=0$.
\end{itemize}
Di conseguenza, il numero di parametri continui liberi passa da $9$ a
\[
9-6=3.
\]
Il vincolo $\det R=1$ seleziona la componente delle rotazioni (escludendo le riflessioni), ma non cambia il conteggio dei gradi di libertà continui.

\medskip
\textbf{Interpretazione geometrica.}
Una rotazione nello spazio è determinata da $3$ parametri (ad esempio tre angoli), quindi $R$ dipende da $3$ gradi di libertà.\\
\fineapprof
Scegliamo la matrice $R(t)$ variabile nel tempo tale che moltiplicata a sinistra per  il vettore $\underline{x_q}(0) - \underline{x_p}(0)$ dia la posizione del punto $q$ rispetto al punto $p$ all'istante $t$, ovvero il vettore $\underline{x_q}(t) - \underline{x_p}(t)$. In formule
\[
\underline{x}_q(t)=\underline{x}_p(t)+R\bigl(\underline{x}_q(0)-\underline{x}_p(0)\bigr).
\]

\begin{oss}{}
    Per avere $\underline{x_q}(t = 0) = \underline{x_q}(0)$ necessariamente si deve avere $R(0) = I$.
\end{oss}

\begin{center}
\begin{tikzpicture}[scale=1.05, line cap=round, line join=round]

% assi in prospettiva
\coordinate (O) at (0,0);
\draw[->,line width=1.1pt] (O) -- (-2.4,-1.4) node[below left] {$x$};
\draw[->,line width=1.1pt] (O) -- (4.6,0) node[below right] {$y$};
\draw[->,line width=1.1pt] (O) -- (0,3.8) node[above] {$z$};

% punti
\coordinate (Xp) at (3.0,1.35);
\coordinate (Xq) at (3.35,2.75);

% vettori posizione
\draw[->,line width=1.2pt] (O) -- (Xp);
\draw[->,line width=1.2pt] (O) -- (Xq);

% vettore differenza x_q - x_p
\draw[->,line width=1.2pt] (Xp) -- (Xq);

% punti marcati
\fill (Xp) circle (2.2pt);
\fill (Xq) circle (2.2pt);

% etichette
\node[below right] at (Xp) {$\underline{x}_p$};
\node[above left]  at (Xq) {$\underline{x}_q$};
\node[right] at ($(Xp)!0.55!(Xq)$) {$\underline{x}_q-\underline{x}_p$};

\end{tikzpicture}
\end{center}

determiniamo la velocità:
\[
(\underline{x}_q-\underline{x}_p)(t)=R\bigl(\underline{x}_q(0)-\underline{x}_p(0)\bigr)
\]
\[
\Rightarrow\quad
R^{\mathsf T}(\underline{x}_q-\underline{x}_p)(t)=\underline{x}_q(0)-\underline{x}_p(0)
\]
\[
\Rightarrow\quad
\frac{d\underline{x}_q}{dt}(t)=\frac{d\underline{x}_p}{dt}+\frac{d}{dt}\Bigl(R\bigl(\underline{x}_q(0)-\underline{x}_p(0)\bigr)\Bigr)
\]
\[
\Rightarrow\quad
\underline{v}_q=\underline{v}_p+\dot{R}R^{\mathsf T}(\underline{x}_q-\underline{x}_p).
\]
(dove $R=(r_{ij})$)

\medskip

$R(t)$ e $R^{\mathsf T}(t)$ sono funzioni del tempo e $RR^{\mathsf T}=I$.
Quindi
\[
\frac{d}{dt}(RR^{\mathsf T})=0
\ \Longleftrightarrow\
\dot{R}R^{\mathsf T}+R\dot{R}^{\mathsf T}=0
\ \Longleftrightarrow\
\dot{R}R^{\mathsf T}=-R\dot{R}^{\mathsf T}=-(\dot{R}R^{\mathsf T})^{\mathsf T}.
\]
\[
\Rightarrow\quad \dot{R}R^{\mathsf T}\ \text{è antisimmetrica, cioè}\quad
\dot{R}R^{\mathsf T}=
\begin{pmatrix}
0 & -\omega_z & \omega_y\\
\omega_z & 0 & -\omega_x\\
-\omega_y & \omega_x & 0
\end{pmatrix}.
\]

Quindi la velocità si può scrivere come
\[
\underline{v}_q=\underline{v}_p+
\begin{pmatrix}
0 & -\omega_z & \omega_y\\
\omega_z & 0 & -\omega_x\\
-\omega_y & \omega_x & 0
\end{pmatrix}
\begin{pmatrix}
x_q-x_p\\
y_q-y_p\\
z_q-z_p
\end{pmatrix}
=
\underline{v}_p+
\begin{pmatrix}
\omega_y(z_q-z_p)-\omega_z(y_q-y_p)\\
\omega_z(x_q-x_p)-\omega_x(z_q-z_p)\\
\omega_x(y_q-y_p)-\omega_y(x_q-x_p)
\end{pmatrix}.
\]

\[
\underline{v}_q=\underline{v}_p+\underline{\omega}\wedge(\underline{x}_q-\underline{x}_p),
\qquad
\underline{\omega}=
\begin{pmatrix}
\omega_x\\
\omega_y\\
\omega_z
\end{pmatrix}
\quad\text{(velocità angolare).}
\]

\[
\text{\emph{Formula fondamentale del moto dei corpi rigidi}.}
\]

\begin{oss}{}
$\underline{x}_p$ e $\underline{x}_q$ sono due punti qualsiasi del corpo rigido.
\end{oss}

\begin{oss}{}
In generale $\underline{\omega}=\bigl(\omega_x(t),\omega_y(t),\omega_z(t)\bigr)$, ma se vincoliamo il moto nel piano allora
\[
\underline{\omega}=(0,0,\omega(t)).
\]
\end{oss}
Possiamo derivare nel tempo la formula fondamentale ottenendo
\[
\frac{d}{dt}\left(\underline{v_q}\right)
=
\frac{d}{dt}\left(\underline{v_p}\right)
+
\frac{d}{dt}\left(\underline{\omega}\wedge\left(\underline{x_q}-\underline{x_p}\right)\right)
=
\]

\[
=
\underline{a_p}
+
\frac{d}{dt}\left(\underline{\omega}\right)\wedge\left(\underline{x_q}-\underline{x_p}\right)
+
\underline{\omega}\wedge\frac{d}{dt}\left(\underline{x_q}-\underline{x_p}\right)
=
\]

\[
=
\underline{a_p}
+
\underline{\alpha}\wedge\left(\underline{x_q}-\underline{x_p}\right)
+
\underline{\omega}\wedge\left(\underline{v_q}-\underline{v_p}\right)
=
\]

\[
=
\underline{a_p}
+
\underline{\alpha}\wedge\left(\underline{x_q}-\underline{x_p}\right)
+
\underline{\omega}\wedge\left(\underline{v_p}+\underline{\omega}\wedge\left(\underline{x_q}-\underline{x_p}\right)-\underline{v_p}\right)
=
\]

\[
=
\underline{a_p}
+
\underline{\alpha}\wedge\left(\underline{x_q}-\underline{x_p}\right)
+
\underline{\omega}\wedge\left(\underline{\omega}\wedge\left(\underline{x_q}-\underline{x_p}\right)\right)
\]
ovvero

\[
\underline{a_q}
=
\underline{a_p}
+
\underline{\alpha}\wedge\left(\underline{x_q}-\underline{x_p}\right)
+
\underline{\omega}\wedge\left(\underline{\omega}\wedge\left(\underline{x_q}-\underline{x_p}\right)\right)
\]
